% --- Archon physics formalization source begin ---
% archon:physics
% archon:covers IPhO2026Problems/problem_IPhO_2026_3_C_5.lean
% archon:source-report reports/ipho_2026_k3/problem_IPhO_2026_3_C_5.source.json
% archon:problem-id IPhO_2026_3
% archon:part-id C.5
% archon:previous-part IPhO_2026_3_C_4
% archon:previous-part-policy natural_language_prerequisite_only

\chapter{Physics problem IPhO\_2026\_3\_C\_5}
\label{ch:IPhO2026Problems_problem_IPhO_2026_3_C_5}

\paragraph{Problem source.}
The paramagnetic torus executes the Carnot refrigeration cycle
1 -> 2 -> 3 -> 4 -> 1 shown in Figure 3b in the H-versus-T plane.  T\_h and T\_c
are the hot- and cold-reservoir temperatures; Q\_h is the magnitude of heat
delivered to the hot reservoir and Q\_c is the magnitude absorbed from the cold
reservoir.  The equation of state is T*M*V = n*K*H and the isothermal heat
relation from part B may be reused.

Current subquestion:
Determine the overall coefficient of performance COP = Q\_c/W for all cycles up to the time found in C4.

\paragraph{Current subquestion.}
Determine the overall coefficient of performance COP = Q\_c/W for all cycles up to the time found in C4.

\paragraph{Recorded answer/context.}
COP = [(T\_h/(T\_0 - T))*ln(T\_0/T) - 1]\textasciicircum{}(-1).

\paragraph{Figure/image path.}
/root/proposal\_for\_physic/science-mango/ipho\_2026\_source/image/T3\_page-4.png

\paragraph{Reusable previous-part conclusions.}
\begin{itemize}
\item Source C.4. Question: A body of heat capacity C\_c is cooled from T\_0 to T while refrigerator input power P and hot-reservoir temperature T\_h remain constant. Determine the elapsed time. Reusable conclusions: t = (C\_c*T\_h/P)*[ln(T\_0/T) - (T\_0 - T)/T\_h]. Policy: natural\_language\_prerequisite\_only; do\_not\_import\_Lean\_output
\end{itemize}

\paragraph{Formalization target.}
create a compiling Lean file with sorry bodies at `IPhO2026Problems/problem\_IPhO\_2026\_3\_C\_5.lean`.
The Lean declarations must preserve the physical quantities, dimensions or dimensional roles, figure labels, governing-law hypotheses, and final relation expressed by this problem.
Use Mathlib/Physlib names found through LeanExplore where available. If a domain API is missing, introduce faithful local abstractions rather than scalar placeholder aliases.

\begin{theorem}[Physics formalization target]
\label{thm:physics:IPhO_2026_3_C_5:target}
The assigned autoformalize agent should translate this physics problem into Lean declarations in the covered file, with theorem and lemma proof bodies written as `by sorry`.
\end{theorem}
\begin{proof}
This is an autoformalization task, not a proof task. Produce faithful statements that can later be proved without weakening the source contract.
\end{proof}
% --- Archon physics formalization source end ---
